% ============================================================
%  SECCIÓN 5 – HISTORIAS ÉPICAS E HISTORIAS DE USUARIO
% ============================================================
\section{Historias Épicas e Historias de Usuario}

En ingeniería de software ágil, una \textbf{historia épica} agrupa un conjunto de
\textbf{historias de usuario} relacionadas que, en conjunto, cubren una capacidad funcional
mayor del sistema. Para el proyecto \textit{!Blood} se identificaron doce historias épicas
que abarcan la totalidad del alcance funcional de la aplicación.

% ─────────────────────────────────────────────────────────────
\subsection{Historias Épicas}

La siguiente tabla presenta las doce historias épicas de \textit{!Blood},
identificadas mediante el prefijo \textbf{HE} seguido de un número secuencial.

\vspace{0.5em}

\begin{longtable}{%
  |>{\bfseries\centering\arraybackslash}p{1.6cm}%
  |>{\raggedright\arraybackslash}p{4.2cm}%
  |>{\raggedright\arraybackslash}p{9.0cm}|}

\caption{Historias épicas del sistema \textit{!Blood}.}
\label{tab:historias_epicas}\\

\hline
\rowcolor{darkgray}
\textcolor{white}{\textbf{ID}} &
\textcolor{white}{\textbf{Historia Épica}} &
\textcolor{white}{\textbf{Descripción}} \\
\hline
\endfirsthead

\hline
\rowcolor{darkgray}
\textcolor{white}{\textbf{ID}} &
\textcolor{white}{\textbf{Historia Épica}} &
\textcolor{white}{\textbf{Descripción}} \\
\hline
\endhead

\hline
\endfoot

\hline
\endlastfoot

HE-01 &
Gestión de Acceso de Usuarios &
El sistema debe permitir gestionar el acceso de los usuarios a la plataforma
\textit{!Blood}: registro en tres pasos con selector de tipo de documento, fotografías del
documento y selfie con cámara, inicio de sesión y recuperación de contraseña. \\
\hline

HE-02 &
Gestión del Perfil de Usuario &
El sistema debe permitir a cada usuario consultar y actualizar su información personal,
tipo de sangre, calificación por estrellas, estadísticas de donación, encuesta de aptitud
y configuración de cuenta. \\
\hline

HE-03 &
Gestión de Solicitudes de Sangre &
El sistema debe permitir crear, publicar, explorar y gestionar solicitudes de donación de
sangre con filtros por compatibilidad y cercanía geográfica. El detalle de la solicitud
muestra el botón \textit{Quiero Donar} siempre visible sin necesidad de desplazamiento. \\
\hline

HE-04 &
Gestión del Proceso de Donación &
El sistema debe guiar al donante en la confirmación de donación (lista de verificación de
salud) e informar que la cancelación posterior implica una penalización de 30 días sin
poder donar. \\
\hline

HE-05 &
Sistema de Calificación Mutua &
El sistema debe permitir que, tras completarse una donación, tanto el donante como el
solicitante se califiquen mutuamente con 1--5 estrellas, un aviso de equidad y un
comentario opcional, contribuyendo a la reputación visible en el perfil. \\
\hline

HE-06 &
Gestión de Campañas de Donación &
El sistema debe permitir a los profesionales de salud crear, publicar y gestionar campañas
institucionales de donación, y a todos los usuarios explorarlas y reservar cupos. Incluye
las vistas exclusivas del profesional: \textit{Home Pro}, \textit{Mis Campañas},
\textit{Gestión de Campaña} y \textit{Perfil Pro}. \\
\hline

HE-07 &
Sistema de Notificaciones &
El sistema debe notificar en tiempo real sobre solicitudes urgentes compatibles, aceptación
de donación por parte de un donante, recordatorio de campaña próxima y reconocimientos de
logros. \\
\hline

HE-08 &
Moderación y Seguridad de Contenido &
El sistema debe garantizar la privacidad de los datos de los usuarios, gestionar el proceso
de verificación de identidad (foto frontal y trasera del documento más selfie) y prevenir
el uso indebido de la plataforma. \\
\hline

HE-09 &
Estadísticas y Gamificación &
El sistema debe mostrar estadísticas individuales y comunitarias, otorgar insignias,
mostrar la calificación promedio por estrellas en el perfil y enviar mensajes
motivacionales para fomentar la donación recurrente. \\
\hline

HE-10 &
Compatibilidad Sanguínea y Geolocalización &
El sistema debe determinar automáticamente la compatibilidad entre tipos de sangre y
calcular la distancia entre el usuario y cada solicitud o campaña. \\
\hline

HE-11 &
Sistema de Penalizaciones &
El sistema debe aplicar penalizaciones automáticas por: \textit{(1)} abuso de urgencia en
solicitudes, que resulta en bloqueo temporal de creación de solicitudes; \textit{(2)}
cancelación de donación confirmada, que impide confirmar donaciones durante 30 días; y
\textit{(3)} calificación discriminatoria o falsa, que conlleva la suspensión de la cuenta. \\
\hline

HE-12 &
Encuesta de Aptitud para Donación &
El sistema debe presentar al usuario una encuesta de ocho preguntas de sí/no tras completar
el registro (Paso~3), determinar su aptitud para donar y permitirle retomar la encuesta
desde su perfil en cualquier momento. \\
\hline

\end{longtable}

% ─────────────────────────────────────────────────────────────
\subsection{Historias de Usuario}

Las historias de usuario detallan, para cada historia épica, los criterios de aceptación
mediante escenarios de prueba funcional (formato \textit{Dado -- Cuando -- Entonces}).
El documento completo, que incluye las doce épicas con sus respectivas historias de usuario
y criterios de aceptación, está disponible en el siguiente enlace:

\vspace{0.4em}
\begin{center}
  \href{https://drive.google.com/uc?id=1Rqvc8YO-pXUc7-rpU93Pk1LDsg7gfyI4&export=download}{%
    \textcolor{bloodred}{\textbf{Descargar Excel de Historias Épicas e Historias de Usuario}}%
  }
\end{center}
\vspace{0.4em}

El archivo Excel contiene una hoja por historia épica, con columnas de identificador,
rol del usuario, enunciado (\textit{Quiero}), propósito (\textit{Para}) y escenarios de
criterios de aceptación.
